\subsection{Ключевые алгоритмы и структуры данных}

\subsubsection{Алгоритм работы классического пула потоков}

Классический пул потоков построен на трёх взаимодействующих
компонентах: очереди задач, защищённой мьютексом, переменной
условия для уведомления рабочих потоков и атомарном флаге
остановки. Рассмотрим алгоритмы основных операций.

Алгоритм постановки задачи в пул (метод \texttt{submit})
представлен в листинге~\ref{alg:simple-submit}.

\begin{lstlisting}[language=C++,
                   caption={Алгоритм постановки задачи
                            в SimpleThreadPool},
                   label={alg:simple-submit}]
template<typename F, typename... Args>
auto submit(F&& f, Args&&... args)
    -> std::future<invoke_result_t<F, Args...>>
{
    using return_type = invoke_result_t<F, Args...>;

    auto task = make_shared<packaged_task<return_type()>>(
        [f = forward<F>(f),
         args = make_tuple(forward<Args>(args)...)]
        () mutable {
            return apply(move(f), move(args));
        });

    future<return_type> result = task->get_future();

    {
        lock_guard<mutex> lock(queue_mutex_);

        if (stop_.load()) {
            throw runtime_error("Pool is stopped");
        }

        task_queue_.emplace([task]() { (*task)(); });
    }

    cv_.notify_one();
    return result;
}
\end{lstlisting}

Алгоритм рабочего цикла потока состоит из следующих шагов:

\begin{enumerate}
    \item Захватить мьютекс очереди задач.
    \item Ждать на переменной условия пока не выполнится
          предикат: очередь непуста
          ($\neg\,queue.empty()$) или получен сигнал
          остановки ($stop = \text{true}$).
    \item Если получен сигнал остановки и очередь пуста ---
          завершить работу потока.
    \item Извлечь задачу из начала очереди (FIFO).
    \item Освободить мьютекс.
    \item Выполнить задачу.
    \item Перейти к шагу~1.
\end{enumerate}

Реализация рабочего цикла на языке C++ представлена
в листинге~\ref{alg:simple-worker-cpp}.

\begin{lstlisting}[language=C++,
                   caption={Рабочий цикл потока SimpleThreadPool},
                   label={alg:simple-worker-cpp}]
void SimpleThreadPool::worker_thread()
{
    while (true) {
        function<void()> task;

        {
            unique_lock<mutex> lock(queue_mutex_);

            cv_.wait(lock, [this] {
                return stop_.load() || !task_queue_.empty();
            });

            if (stop_.load() && task_queue_.empty()) {
                return;
            }

            task = move(task_queue_.front());
            task_queue_.pop();
        }

        task();
    }
}
\end{lstlisting}

Отметим ключевые свойства данного алгоритма.

Первое свойство --- задача извлекается из очереди под защитой
мьютекса, а выполняется вне его. Это принципиально важно: если
бы задача выполнялась под мьютексом, другие потоки не могли бы
одновременно извлекать задачи из очереди, что свело бы к нулю
преимущества параллельного выполнения.

Второе свойство --- использование \texttt{cv\_.wait()} с предикатом
защищает от ложных пробуждений. Стандарт C++ допускает что
\texttt{condition\_variable} может пробудить поток без
соответствующего уведомления~\cite{Williams2019}. Предикат
гарантирует что поток вернётся в состояние ожидания если условие
не выполнено.

Третье свойство --- порядок извлечения задач FIFO. Задачи
выполняются в том порядке, в котором были поставлены. Это
обеспечивает предсказуемое поведение, но не позволяет
адаптироваться к неравномерной нагрузке.

Алгоритм корректного завершения работы пула состоит
из трёх шагов:

\begin{enumerate}
    \item Захватить мьютекс и установить флаг остановки
          \texttt{stop\_ = true}. Захват мьютекса гарантирует
          что рабочий поток либо уже ожидает на переменной
          условия, либо увидит флаг при следующем захвате
          мьютекса.
    \item Вызвать \texttt{cv\_.notify\_all()} для пробуждения
          всех ожидающих потоков.
    \item Вызвать \texttt{join()} для каждого рабочего потока,
          ожидая его завершения.
\end{enumerate}

Реализация деструктора представлена
в листинге~\ref{alg:simple-destructor}.

\begin{lstlisting}[language=C++,
                   caption={Алгоритм завершения работы
                            SimpleThreadPool},
                   label={alg:simple-destructor}]
SimpleThreadPool::~SimpleThreadPool()
{
    {
        lock_guard<mutex> lock(queue_mutex_);
        stop_.store(true);
    }

    cv_.notify_all();

    for (thread& worker : workers_) {
        if (worker.joinable()) {
            worker.join();
        }
    }
}
\end{lstlisting}

Важным аспектом алгоритма завершения является установка флага
\texttt{stop\_} под защитой мьютекса. Данное решение устраняет
гонку данных между деструктором и рабочими потоками. Без захвата
мьютекса возможна следующая ситуация: деструктор устанавливает
флаг и вызывает \texttt{notify\_all()}, но рабочий поток ещё не
вошёл в \texttt{cv\_.wait()} --- в результате уведомление теряется
и поток блокируется навсегда. Захват мьютекса гарантирует что
в момент установки флага поток либо уже ожидает на переменной
условия (и будет разбужен), либо ещё не захватил мьютекс
(и увидит флаг сразу после его захвата)~\cite{Williams2019}.

Рабочие потоки при получении сигнала остановки продолжают
выполнять задачи из очереди до её полного опустошения, после
чего завершаются. Это обеспечивает корректное завершение
(graceful shutdown): все задачи, поставленные до вызова
деструктора, будут выполнены.